%%%%%%%%%%  scan idx 180  |  folio 167  %%%%%%%%%%
$0\to L_1\to L_0\to L\to0$. D'%
après ce qu'on a vu plus haut, on a de même
\[
F^n(L)=\varinjlim\operatorname{Coker}(F^{n-1}(L_0)\to F^{n-1}(L_1)).
\]
comme $\varphi=\varphi^0:F^0\to G^0$ est un isomorphisme, il en résulte, par récurrence sur $n$, que le morphisme $\varphi^n:F^n\to G^n$ induit par $\varphi$ est un isomorphisme pour tout $n$.

b) L'équivalence (i)$\Leftrightarrow$(ii) résulte d'un dévissage immédiat, laissé au lecteur.

\textbf{Remarques 2.1.3.1.} On a bien entendu une équivalence duale de a), qu'on laisse au lecteur le soin d'énoncer. D'autre part, si $\mathcal A$ et $\mathcal B$ ont assez d'injectifs, les conditions (i) et (ii) de b) sont équivalentes à
\begin{enumerate}[label=(\roman*),start=3]
\item pour $L\in\Ob\mathrm D^{-}(\mathcal A)$ et $M\in\Ob\mathrm D^{+}(\mathcal A)$, la flèche canonique
\[
\mathbb R\!\operatorname{Hom}(L,M)\longrightarrow
\mathbb R\!\operatorname{Hom}(\varphi(L),\varphi(M))
\]
est un isomorphisme,
\end{enumerate}
comme on le voit aussitôt par un argument de suite spectrale.

\subsection*{2.2. Schémas noethériens et schémas divisoriels}

\textbf{Notations 2.2.1.} Soit $S$ un schéma. On note $\Qcoh(S)$ (resp. $\Coh(S)$) la catégorie (abélienne) des $\mathcal O_S$-Modules quasi-cohérents (resp. cohérents). Soit $[m,n]$ un intervalle de $\mathbb Z$. On dira qu'un objet $E$ de $\mathrm D(\Qcoh(S))$ est $n$-pseudo-cohérent (resp. pseudo-cohérent, resp. d'amplitude parfaite contenue dans $[m,n]$, resp. parfait) si $E$ est $n$-pseudo-cohérent (resp. pseudo-cohérent, resp. d'amplitude parfaite contenue dans $[m,n]$, resp. parfait) relativement à la catégorie fibrée sur le site zariskien de $S$ définie par $U\mapsto\Loclib(U)$ ((I 2.3), (I 4.7) et (2.0)). On désignera par $\mathrm D^{-}(\Qcoh(S))_{\coh}$ (resp. $\mathrm D(\Qcoh(S))_{\parf}$) la sous-catégorie triangulée pleine de $\mathrm D(\Qcoh(S))$ formée des complexes pseudo-cohérents (resp. parfaits).

\textbf{Proposition 2.2.2}\footnote{Ce résultat, accessoire, ne sera pas utilisé dans la suite du Séminaire : on recommande au lecteur de le négliger en première lecture.}. \emph{Soit $S$ un schéma noethérien. Le foncteur}
\emph{canonique}
